\chapter{Introduction}
\label{ch:introduction}

Black-box machine learning models are increasingly being integrated
into high-impact real-world settings, with healthcare being a prominent example.
These technologies hold the potential to enhance diagnostic accuracy,
improve workflow efficiency, and support clinical decision-making.
However, in such safety-critical environments,
it is essential not only to evaluate predictive performance
but also to quantify model uncertainty in a meaningful way.

Standard evaluation metrics often fall short
by primarily focusing on point predictions
and lacking the ability to express the model's confidence in its outputs.
This lack of calibrated uncertainty can be especially problematic in domains where
incorrect decisions carry severe consequences for patient health.
For example, in medical imaging tasks involving multi-label classification,
a single image may present multiple conditions,
increasing the complexity of delivering reliable assessments.
In such scenarios, understanding and managing uncertainty
is essential for building trustworthy AI systems.

Addressing the critical need for dependable confidence measures
falls under the domain of Uncertainty Quantification (UQ).
Within this field, Conformal Prediction (CP) stands out as a particularly appealing framework.
It offers a mathematically rigorous yet highly flexible approach
that can work with virtually any underlying machine learning model.
Rather than committing to a single prediction,
CP generates prediction sets-a range or set of possible outcomes-
accompanied by a guarantee that the true outcome will be contained
within the set with a user-specified probability (e.g., 90\%).
This model-agnostic nature makes CP especially useful for
enhancing the reliability of complex black-box systems,
without requiring changes to the underlying architecture.

While the coverage guarantees provided by standard CP are broadly useful,
critical applications like medical diagnosis
often demand more nuanced control over potential errors.
In many clinical situations, certain types of mistakes-
such as failing to detect an existing condition (a false negative)-
are far more detrimental than others.
Conformal Risk Control (CRC) extends the core ideas of CP to address this need.
Rather than providing blanket coverage guarantees,
CRC enables practitioners to control the average rate of specific,
user-defined errors that align with domain-specific priorities.
This ability to tailor uncertainty quantification to clinically meaningful risk types
makes CRC a particularly compelling solution for safety-critical settings.

Recognizing the potential of CRC to provide more meaningful safety assurances in clinical settings,
this project investigates the application of Conformal Risk Control
to multi-label classification tasks within medical imaging.
The goal is to explore its effectiveness in managing diagnostic risks
in a way that aligns more closely with clinical needs than traditional
evaluation metrics or standard uncertainty quantification approaches.
